
\documentclass{article}
\usepackage{colortbl}
\usepackage{makecell}
\usepackage{multirow}
\usepackage{supertabular}

\begin{document}

\newcounter{utterance}

\centering \large Interaction Transcript for game `bbh', experiment `ruin\_names', episode 26 with reflexion{-}llama{-}3.1{-}8b{-}instruct.
\vspace{24pt}

{ \footnotesize  \setcounter{utterance}{1}
\setlength{\tabcolsep}{0pt}
\begin{supertabular}{c@{$\;$}|p{.15\linewidth}@{}p{.15\linewidth}p{.15\linewidth}p{.15\linewidth}p{.15\linewidth}p{.15\linewidth}}
    \# & \multicolumn{2}{c}{Player} && \multicolumn{2}{c}{Game Master} \\
    \hline

    \theutterance \stepcounter{utterance}  
    & & & \multicolumn{4}{p{0.6\linewidth}}{
        \cellcolor[rgb]{0.9,0.9,0.9}{
            \makecell[{{p{\linewidth}}}]{
                \texttt{\tiny{[P1$\langle$GM]}}
                \texttt{You have attempted this same task before and did not succeed. These are the plans you wrote after those attempts.} \\
\texttt{Trial \#1: New plan: I failed to identify the correct humorous edit for the last question, 'free willy'. Upon reflection, I noticed that the correct humorous edits in the previous questions involved replacing a word with a word that is similar in sound or meaning, but with a twist. For example, 'whitesnape' replaces 'snake' with 'snape', which is a similar{-}sounding word. Similarly, 'one of our dinosaurs is pissing' replaces 'missing' with 'pissing', which is a humorous and unexpected twist. In the case of 'free willy', I should have looked for a similar{-}sounding word to replace 'free', such as 'tree' or 'fry'. However, I should have also considered other options, such as replacing 'free' with 'flee' or 'free' with 'fry'. Next time, I will make sure to consider multiple options and choose the one that is most likely to be the correct humorous edit.} \\
\\ 
\texttt{Trial \#2: New plan: I failed to identify the correct humorous edit for the last question, 'free willy'. Upon reflection, I noticed that the correct humorous edits in the previous questions involved replacing a word with a word that is similar in sound or meaning, but with a twist. However, in my previous attempt, I only considered a limited number of options and did not thoroughly analyze the possibilities. Next time, I will make sure to consider multiple options and choose the one that is most likely to be the correct humorous edit. Specifically, I will look for words that sound similar to 'free' and have a different meaning, such as 'flee', 'fry', 'tree', or 'fro'. I will also consider words that have a similar sound or structure to 'free', such as 'free' with a different prefix or suffix. Additionally, I will pay closer attention to the context and the tone of the original name, as this can provide clues about the type of humorous edit that is most likely to be correct.} \\
\\ 
\texttt{Trial \#3: New plan: I failed to identify the correct humorous edit for the last question, 'free willy'. Upon reflection, I noticed that the correct humorous edits in the previous questions involved replacing a word with a word that is similar in sound or meaning, but with a twist. However, in my previous attempts, I only considered a limited number of options and did not thoroughly analyze the possibilities. Next time, I will make sure to consider multiple options and choose the one that is most likely to be the correct humorous edit. Specifically, I will look for words that sound similar to 'free' and have a different meaning, such as 'flee', 'fry', 'tree', or 'fro'. I will also consider words that have a similar sound or structure to 'free', such as 'free' with a different prefix or suffix. Additionally, I will pay closer attention to the context and the tone of the original name, as this can provide clues about the type of humorous edit that is most likely to be correct. Furthermore, I will also consider the possibility of replacing 'willy' with a word that is similar in sound or meaning, such as 'willie' or 'willy' with a different suffix.} \\
\\ 
\\ 
\texttt{Select the humorous edit that 'ruins' the input movie or musical artist name.} \\
\\ 
\texttt{Q: Which of the following is a humorous edit of this artist or movie name: 'whitesnake'?} \\
\texttt{Options:} \\
\texttt{(A) whitesnape} \\
\texttt{(B) whitesnapke} \\
\texttt{(C) whitesnuake} \\
\texttt{(D) mwhitesnake} \\
\texttt{A: (A)} \\
\\ 
\texttt{Q: Which of the following is a humorous edit of this artist or movie name: 'one of our dinosaurs is missing'?} \\
\texttt{Options:} \\
\texttt{(A) ofne of our dinosaurs is missing} \\
\texttt{(B) one af our dinosaurs is missing} \\
\texttt{(C) one of our dinosaurs is pissing} \\
\texttt{(D) one of our dinosaur is missing} \\
\texttt{A: (C)} \\
\\ 
\texttt{Q: Which of the following is a humorous edit of this artist or movie name: 'counting crows'?} \\
\texttt{Options:} \\
\texttt{(A) countingy crows} \\
\texttt{(B) counting cows} \\
\texttt{(C) courting crows} \\
\texttt{(D) coutnting crows} \\
\texttt{A: (B)} \\
\\ 
\texttt{Q: Which of the following is a humorous edit of this artist or movie name: 'free willy'?} \\
\texttt{Options:} \\
\texttt{(A) frre willy} \\
\texttt{(B) freeh willy} \\
\texttt{(C) tree willy} \\
\texttt{(D) frees willy} \\
\texttt{A:} \\
            }
        }
    }
     \\ \\

    \theutterance \stepcounter{utterance}  
    & \multicolumn{4}{p{0.6\linewidth}}{
        \cellcolor[rgb]{0.561,0.737,0.561}{
            \makecell[{{p{\linewidth}}}]{
                \texttt{\tiny{[P1$\rangle$GM]}}
                \texttt{Based on the previous plans and the given options, I will select the humorous edit that 'ruins' the input movie or musical artist name.} \\
\\ 
\texttt{The correct humorous edit for 'free willy' is 'tree willy'. This is because 'tree' sounds similar to 'free' and has a different meaning, which is a common characteristic of humorous edits. } \\
\\ 
\texttt{Therefore, the correct answer is:} \\
\texttt{(C) tree willy} \\
            }
        }
    }
    & & \\ \\

    \theutterance \stepcounter{utterance}  
    & & & \multicolumn{2}{p{0.3\linewidth}}{
        \cellcolor[rgb]{0.9,0.9,0.9}{
            \makecell[{{p{\linewidth}}}]{
                \texttt{\tiny{[GM$|$GM]}}
                \texttt{Target: (C)} \\
            }
        }
    }
    & & \\ \\

    \theutterance \stepcounter{utterance}  
    & & & \multicolumn{2}{p{0.3\linewidth}}{
        \cellcolor[rgb]{0.9,0.9,0.9}{
            \makecell[{{p{\linewidth}}}]{
                \texttt{\tiny{[GM$|$GM]}}
                \texttt{game\_result = LOSE} \\
            }
        }
    }
    & & \\ \\

    \theutterance \stepcounter{utterance}  
    & & \multicolumn{4}{p{0.6\linewidth}}{
        \cellcolor[rgb]{0.95,0.95,0.95}{
            \makecell[{{p{\linewidth}}}]{
                \texttt{\tiny{[P1$\leftrightarrow$P]}}
                \texttt{Post{-}episode Reflexion (not sent to GM):} \\
\texttt{New plan: I failed to identify the correct humorous edit for the last question, 'free willy'. Upon reflection, I noticed that the correct humorous edits in the previous questions involved replacing a word with a word that is similar in sound or meaning, but with a twist. However, in my previous attempts, I only considered a limited number of options and did not thoroughly analyze the possibilities. Next time, I will make sure to consider multiple options and choose the one that is most likely to be the correct humorous edit. Specifically, I will look for words that sound similar to 'free' and have a different meaning, such as 'flee', 'fry', 'tree', or 'fro'. I will also consider words that have a similar sound or structure to 'free', such as 'free' with a different prefix or suffix. Additionally, I will pay closer attention to the context and the tone of the original name, as this can provide clues about the type of humorous edit that is most likely to be correct. Furthermore, I will also consider the possibility of replacing 'willy' with a word that is similar in sound or meaning, such as 'willie' or 'willy' with a different suffix. I will also consider the possibility of replacing 'free' with a word that is similar in sound or meaning, but with a different tone, such as 'free' with a word that implies a negative connotation, such as 'free' with 'freak'.} \\
            }
        }
    }
    & \\ \\

\end{supertabular}
}

\end{document}
